\chapter{Conclusões}
\label{cap6}

\section{Conclusões}
\label{sec:conclusao}

\paragraph{}Este trabalho apresentou o desenvolvimento e a validação do \textit{MinervaMesh}, uma plataforma \textit{client-side} para simulação numérica em Engenharia Mecânica, executada integralmente no navegador por meio da combinação de \textit{PyScript}, \textit{Pyodide} e \textit{WebAssembly}. Os objetivos estabelecidos na Introdução foram alcançados: implementaram-se oito módulos de Método dos Elementos Finitos descritos no Capítulo~\ref{cap4}, dois módulos complementares de Método das Diferenças Finitas (Seção~\ref{sec:modulos_mdf}) e uma interface interativa de pré/pós-processamento comum a todas as simulações, validada prioritariamente contra soluções analíticas e, para os casos 2D sem solução fechada, contra \textit{benchmark} consagrado da literatura. Os dois resultados centrais do trabalho confirmam-se: as equações discretizadas pelo MEF e pelo MDF \textbf{reproduzem corretamente as soluções de referência} em todos os casos verificados, e a \textbf{implementação na web flexibiliza a resolução desses problemas para outras pessoas} --- estudantes, monitores e pesquisadores ---, que podem parametrizar, executar e inspecionar cada simulação diretamente no navegador, sem instalação nem custo de licenciamento.

\paragraph{}Os resultados de validação apresentados no Capítulo~\ref{cap5} demonstram que a plataforma preserva a ordem de convergência esperada para cada discretização: erros nodais da ordem da precisão de máquina ($\sim 10^{-13}$) para problemas com solução polinomial --- condução 1D e viga de Euler-Bernoulli --- como consequência da superconvergência de elementos lineares e da reprodução polinomial dos elementos cúbicos de Hermite; concordância visual com soluções analíticas em problemas parabólicos e hiperbólicos; e, para o escoamento em cavidade com tampa móvel a $\text{Re}=100$, erro relativo inferior a 0,4\% na função de corrente mínima $\psi_{\min}$ comparada à referência de Ghia et al.\ \cite{Ghia1982}, utilizando apenas malha $41 \times 41$. A estabilização SUPG reproduz quantitativamente o parâmetro $\tau$ canônico da formulação de Brooks e Hughes \cite{Brooks1982}, e a formulação Vorticidade-Corrente com a condição de contorno de Thom \cite{Thom1933} na parede reproduz o vórtice primário da cavidade dentro da resolução da malha utilizada.

\paragraph{}A contribuição didática do trabalho articula-se em três eixos: (i) \textbf{acessibilidade} --- a plataforma é gratuita, opera sem instalação e pode ser hospedada em serviços de páginas estáticas sem custo operacional; (ii) \textbf{transparência} --- o código-fonte Python de cada simulação é inspecionável pelo usuário diretamente na página, permitindo o uso tanto como ferramenta de simulação quanto como objeto de estudo em métodos numéricos; e (iii) \textbf{unificação} --- os três domínios clássicos abordados (térmico, fluidodinâmico e estrutural) são cobertos em uma única interface consistente, facilitando comparações pedagógicas entre métodos e formulações para um mesmo fenômeno físico. A esses três eixos didáticos soma-se uma contribuição metodológica: o \textit{benchmark} cruzado entre \textit{Pyodide} no navegador e \textit{CPython} nativo apresentado na Seção~\ref{sec:analise_computacional} caracteriza quantitativamente o custo da execução \textit{client-side} e oferece um protocolo reproduzível para análises semelhantes em outras plataformas didáticas baseadas em \textit{PyScript}/\textit{Pyodide}.

\paragraph{}A principal limitação da arquitetura \textit{client-side} é a capacidade computacional: embora o \textit{WebAssembly} opere em velocidade próxima à nativa para as rotinas de \textit{NumPy} e \textit{SciPy} utilizadas, malhas com dezenas de milhares de nós impõem tempos de resposta incompatíveis com a interatividade esperada em ambiente didático. O \textit{benchmark} cruzado entre \textit{Pyodide} e \textit{CPython} nativo apresentado na Seção~\ref{sec:analise_computacional} confirma quantitativamente essa restrição, e a Seção~\ref{sec:limitacoes} sistematiza as fronteiras práticas da plataforma. Para os casos deste trabalho (centenas a poucos milhares de nós), essa restrição não se mostrou limitante, mas condiciona a estratégia de extensão discutida a seguir.

\section{Trabalhos Futuros}

\paragraph{}A partir da base estabelecida, identificam-se sete linhas naturais de continuidade:

\begin{itemize}
    \item \textbf{Análise de convergência quantitativa:} produzir gráficos log-log de erro $L^2$ em função do refinamento de malha ($h$) para os casos 2D, validando numericamente a ordem de convergência teórica discutida na fundamentação.

    \item \textbf{Elementos de ordem superior:} estender a biblioteca de funções de forma para elementos quadráticos ($\mathbb{P}_2$) e quadriláteros bilineares ($\mathbb{Q}_1$), permitindo estudos comparativos entre discretizações distintas para o mesmo problema físico.

    \item \textbf{Modelos de turbulência:} incorporar formulações RANS ou de \textit{Large Eddy Simulation} ao solver de Navier-Stokes 2D, ampliando o alcance da plataforma para regimes não laminares.

    \item \textbf{Escalabilidade tridimensional:} investigar a viabilidade de solvers 3D sob a restrição de memória e tempo do \textit{WebAssembly}, possivelmente explorando execução assíncrona via \textit{Web Workers} e solvers iterativos esparsos (cf.\ Seção~\ref{sec:limitacoes}).

    \item \textbf{Validação experimental:} complementar o atual processo de verificação numérica com comparações contra dados experimentais, introduzindo a nomenclatura formal de V\&V (Verificação e Validação) e a análise de incerteza na discussão metodológica.

    \item \textbf{Integração com geradores de malha externos:} ampliar o repertório de geometrias utilizáveis incorporando à plataforma o \texttt{triangle-wasm} \cite{TriangleWasm} --- \textit{port} \textit{WebAssembly} do \textit{Triangle} já disponível e funcional --- e, à medida que se viabilize, um eventual \textit{port} do \texttt{Gmsh}, atualmente inexistente em forma oficial. Tal integração endereçaria simultaneamente as restrições de geometria complexa, refinamento adaptativo e extensão tridimensional discutidas na Seção~\ref{sec:limitacoes}, preservando a característica \textit{client-side} da plataforma.

    \item \textbf{Condições de contorno de Robin:} incorporar a imposição de condições mistas (Seção~\ref{sec:condicoes_contorno}) ao módulo de transferência de calor 2D, permitindo a representação de troca convectiva $-k\,\partial T/\partial n = h(T - T_\infty)$ entre o domínio sólido e um meio externo a temperatura prescrita --- extensão natural pedagogicamente relevante para problemas de resfriamento de aletas e dispositivos eletrônicos.
\end{itemize}
